\small
\begin{enumerate}[label=\textbf{\arabic*.}, leftmargin=0pt, itemsep=1.2em]
\item К какому типу алгоритмов машинного обучения относится сеть Кохонена?
\begin{itemize}[label=$\square$, leftmargin=1.5cm, itemsep=0.2em, topsep=0.3em]
\item Обучение с учителем (Supervised Learning)
\item Обучение с подкреплением (Reinforcement Learning)
\item Обучение без учителя (Unsupervised Learning)
\item Полуконтролируемое обучение (Semi-supervised Learning)
\end{itemize}
{\small\textit{Правильный ответ: Обучение без учителя (Unsupervised Learning).}}

\item Как расшифровывается аббревиатура SOM, часто используемая для обозначения сети Кохонена?
\begin{itemize}[label=$\square$, leftmargin=1.5cm, itemsep=0.2em, topsep=0.3em]
\item Standard Output Matrix
\item Self-Organizing Map (Самоорганизующаяся карта)
\item Supervised Optimal Model
\item Statistical Organization Method
\end{itemize}
{\small\textit{Правильный ответ: Self-Organizing Map (Самоорганизующаяся карта).}}

\item Для какой основной задачи применяются самоорганизующиеся карты Кохонена (SOM)?
\begin{itemize}[label=$\square$, leftmargin=1.5cm, itemsep=0.2em, topsep=0.3em]
\item Классификация изображений по заранее известным меткам
\item Аппроксимация нелинейных математических функций
\item Кластеризация данных и проецирование многомерного пространства на пространство меньшей размерности (обычно 2D)
\item Регрессионный анализ временных рядов
\end{itemize}
{\small\textit{Правильный ответ: Кластеризация данных и проецирование многомерного пространства на пространство меньшей размерности (обычно 2D).}}

\item Какая топология нейронов чаще всего используется в картах Кохонена?
\begin{itemize}[label=$\square$, leftmargin=1.5cm, itemsep=0.2em, topsep=0.3em]
\item Одномерная линия
\item Трехмерный куб
\item Двумерная решетка (гексагональная или прямоугольная)
\item Полносвязный ненаправленный граф
\end{itemize}
{\small\textit{Правильный ответ: Двумерная решетка (гексагональная или прямоугольная).}}

\item Какой принцип лежит в основе обучения сети Кохонена?
\begin{itemize}[label=$\square$, leftmargin=1.5cm, itemsep=0.2em, topsep=0.3em]
\item Минимизация функции потерь через градиентный спуск
\item Конкурентное обучение по принципу «победитель получает все» (Winner-Takes-All)
\item Обучение на основе вознаграждения
\item Эволюционный поиск
\end{itemize}
{\small\textit{Правильный ответ: Конкурентное обучение по принципу «победитель получает все» (Winner-Takes-All).}}

\item Что такое BMU (Best Matching Unit) в сети Кохонена?
\begin{itemize}[label=$\square$, leftmargin=1.5cm, itemsep=0.2em, topsep=0.3em]
\item Нейрон, который дальше всего от входного вектора
\item Нейрон-победитель, весовой вектор которого наиболее похож (имеет минимальное расстояние) на текущий входной пример
\item Самый быстро вычисляемый нейрон
\item Нейрон с максимальным значением весов
\end{itemize}
{\small\textit{Правильный ответ: Нейрон-победитель, весовой вектор которого наиболее похож (имеет минимальное расстояние) на текущий входной пример.}}

\item Какая метрика расстояния чаще всего используется для поиска нейрона-победителя?
\begin{itemize}[label=$\square$, leftmargin=1.5cm, itemsep=0.2em, topsep=0.3em]
\item Расстояние Махаланобиса
\item Евклидово расстояние
\item Манхэттенское расстояние
\item Расстояние Чебышева
\end{itemize}
{\small\textit{Правильный ответ: Евклидово расстояние.}}

\item Как происходит адаптация весов в классическом алгоритме обучения SOM?
\begin{itemize}[label=$\square$, leftmargin=1.5cm, itemsep=0.2em, topsep=0.3em]
\item Обновляются веса всех нейронов одинаково
\item Обновляются случайные нейроны
\item Обновляются веса только нейрона-победителя и его топологических соседей, они "подтягиваются" к входному вектору
\item Веса обновляются только если сеть ошибается
\end{itemize}
{\small\textit{Правильный ответ: Обновляются веса только нейрона-победителя и его топологических соседей, они "подтягиваются" к входному вектору.}}

\item Для чего в алгоритме Кохонена используется «функция соседства» (neighborhood function)?
\begin{itemize}[label=$\square$, leftmargin=1.5cm, itemsep=0.2em, topsep=0.3em]
\item Чтобы определить силу изменения весов в зависимости от топологического расстояния нейрона от победителя
\item Для сортировки входных данных
\item Чтобы удалять выбросы из выборки
\item Для ограничения максимального веса
\end{itemize}
{\small\textit{Правильный ответ: Чтобы определить силу изменения весов в зависимости от топологического расстояния нейрона от победителя.}}

\item Какую форму чаще всего имеет функция соседства в сети Кохонена?
\begin{itemize}[label=$\square$, leftmargin=1.5cm, itemsep=0.2em, topsep=0.3em]
\item Линейную
\item Гауссовскую (колокол)
\item Ступенчатую
\item Синусоидальную
\end{itemize}
{\small\textit{Правильный ответ: Гауссовскую (колокол).}}

\item Как меняется радиус окрестности (neighborhood radius) с течением времени (эпох) в ходе обучения SOM?
\begin{itemize}[label=$\square$, leftmargin=1.5cm, itemsep=0.2em, topsep=0.3em]
\item Увеличивается
\item Остается постоянным
\item Изначально большой, затем монотонно убывает (затухает) к нулю
\item Колеблется
\end{itemize}
{\small\textit{Правильный ответ: Изначально большой, затем монотонно убывает (затухает) к нулю.}}

\item Зачем необходимо уменьшать радиус окрестности по мере обучения?
\begin{itemize}[label=$\square$, leftmargin=1.5cm, itemsep=0.2em, topsep=0.3em]
\item Чтобы ускорить выполнение программы
\item Чтобы на начальных этапах происходило грубое глобальное упорядочивание карты, а на поздних — точная локальная подстройка
\item Это требование алгоритма
\item Чтобы предотвратить удаление нейронов
\end{itemize}
{\small\textit{Правильный ответ: Чтобы на начальных этапах происходило грубое глобальное упорядочивание карты, а на поздних — точная локальная подстройка.}}

\item Как скорость обучения (learning rate) обычно ведет себя в процессе обучения карты Кохонена?
\begin{itemize}[label=$\square$, leftmargin=1.5cm, itemsep=0.2em, topsep=0.3em]
\item Она постоянна
\item Она монотонно возрастает
\item Она постепенно затухает, переходя от большого значения к малому
\item Она генерируется случайно на каждой итерации
\end{itemize}
{\small\textit{Правильный ответ: Она постепенно затухает, переходя от большого значения к малому.}}

\item Требуются ли в сети Кохонена скрытые слои (в понимании многослойного персептрона)?
\begin{itemize}[label=$\square$, leftmargin=1.5cm, itemsep=0.2em, topsep=0.3em]
\item Да, обязательно минимум два
\item Нет, архитектура состоит только из входного слоя и слоя конкурентных нейронов (самой карты)
\item Да, количество скрытых слоев должно быть равно размерности карты
\item Сеть состоит только из одного нейрона
\end{itemize}
{\small\textit{Правильный ответ: Нет, архитектура состоит только из входного слоя и слоя конкурентных нейронов (самой карты).}}

\item Нужны ли обучающей выборке для сети Кохонена целевые метки (правильные ответы Y)?
\begin{itemize}[label=$\square$, leftmargin=1.5cm, itemsep=0.2em, topsep=0.3em]
\item Да, иначе градиент будет равен нулю
\item Нет, это алгоритм обучения без учителя (unsupervised)
\item Да, но только для половины данных
\item Это зависит от используемой функции активации
\end{itemize}
{\small\textit{Правильный ответ: Нет, это алгоритм обучения без учителя (unsupervised).}}

\item Если два входных вектора находятся близко друг к другу в исходном многомерном пространстве, где, скорее всего, окажутся их BMU на обученной карте Кохонена?
\begin{itemize}[label=$\square$, leftmargin=1.5cm, itemsep=0.2em, topsep=0.3em]
\item На противоположных концах карты
\item Рядом друг с другом (свойство сохранения топологии)
\item Обязательно в одном и том же нейроне
\item Местоположение в карте будет абсолютно случайным
\end{itemize}
{\small\textit{Правильный ответ: Рядом друг с другом (свойство сохранения топологии).}}

\item Как происходит первоначальная инициализация весов в сети Кохонена?
\begin{itemize}[label=$\square$, leftmargin=1.5cm, itemsep=0.2em, topsep=0.3em]
\item Случайными малыми значениями или выборкой случайных примеров из обучающих данных
\item Все веса устанавливаются в ноль
\item Все веса устанавливаются в единицу
\item Инициализация происходит после первой эпохи
\end{itemize}
{\small\textit{Правильный ответ: Случайными малыми значениями или выборкой случайных примеров из обучающих данных.}}

\item Что произойдет, если размер (сетка) карты Кохонена слишком мал (например, 2x2)?
\begin{itemize}[label=$\square$, leftmargin=1.5cm, itemsep=0.2em, topsep=0.3em]
\item Карта "разорвется"
\item Множество различных кластеров сольются в одни и те же узлы, детализация группирования будет низкой
\item Алгоритм никогда не сойдется
\item Сеть автоматически увеличит свой размер
\end{itemize}
{\small\textit{Правильный ответ: Множество различных кластеров сольются в одни и те же узлы, детализация группирования будет низкой.}}

\item Что означает "мертвый нейрон" в обучении конкурентных сетей?
\begin{itemize}[label=$\square$, leftmargin=1.5cm, itemsep=0.2em, topsep=0.3em]
\item Нейрон, вес которого стал NaN (Not a Number)
\item Нейрон, который физически отключился от RAM
\item Узел сети, который никогда не становится победителем (BMU) и не обновляется
\item Нейрон, выдающий только нули
\end{itemize}
{\small\textit{Правильный ответ: Узел сети, который никогда не становится победителем (BMU) и не обновляется.}}

\item Какой метод инициализации весов помогает избежать проблемы "мертвых нейронов" в SOM?
\begin{itemize}[label=$\square$, leftmargin=1.5cm, itemsep=0.2em, topsep=0.3em]
\item Инициализация нулями
\item Инициализация весов случайными векторами из самой обучающей выборки
\item Инициализация гигантскими значениями
\item Использование отрицательных весов
\end{itemize}
{\small\textit{Правильный ответ: Инициализация весов случайными векторами из самой обучающей выборки.}}

\item Какая визуализация часто используется вместе с сетями Кохонена для отображения границ кластеров?
\begin{itemize}[label=$\square$, leftmargin=1.5cm, itemsep=0.2em, topsep=0.3em]
\item U-матрица (Unified Distance Matrix)
\item Матрица путаницы (Confusion Matrix)
\item ROC-кривая
\item Дендрограмма
\end{itemize}
{\small\textit{Правильный ответ: U-матрица (Unified Distance Matrix).}}

\item Что отображает U-матрица для обученной сети SOM?
\begin{itemize}[label=$\square$, leftmargin=1.5cm, itemsep=0.2em, topsep=0.3em]
\item Веса каждого нейрона
\item Количество примеров, выбравших нейрон победителем
\item Расстояния между весовыми векторами соседних нейронов (чем светлее, тем ближе)
\item Ошибку предсказания на каждом шаге
\end{itemize}
{\small\textit{Правильный ответ: Расстояния между весовыми векторами соседних нейронов (чем светлее, тем ближе).}}

\item Что такое квантизационная ошибка (quantization error) в контексте обученной карты Кохонена?
\begin{itemize}[label=$\square$, leftmargin=1.5cm, itemsep=0.2em, topsep=0.3em]
\item Ошибка округления чисел с плавающей точкой при вычислениях
\item Среднее расстояние между каждым входным вектором и вектором весов его нейрона-победителя (BMU)
\item Количество «мёртвых» нейронов в карте
\item Разность между размером карты и числом кластеров в данных
\end{itemize}
{\small\textit{Правильный ответ: Среднее расстояние между каждым входным вектором и вектором весов его нейрона-победителя (BMU).}}

\item Какая из метрик часто используется для оценки качества разбиения на кластеры без учителя (в т.ч. после использования SOM)?
\begin{itemize}[label=$\square$, leftmargin=1.5cm, itemsep=0.2em, topsep=0.3em]
\item F1-score
\item Silhouette Score (коэффициент силуэта) или индекс Дэвиса-Болдина
\item Accuracy (Точность)
\item R-квадрат
\end{itemize}
{\small\textit{Правильный ответ: Silhouette Score (коэффициент силуэта) или индекс Дэвиса-Болдина.}}

\item С каким популярным алгоритмом кластеризации часто сравнивают сеть Кохонена?
\begin{itemize}[label=$\square$, leftmargin=1.5cm, itemsep=0.2em, topsep=0.3em]
\item Дерево решений
\item k-Means (k-средних)
\item Логистическая регрессия
\item Метод главных компонент (PCA)
\end{itemize}
{\small\textit{Правильный ответ: k-Means (k-средних).}}

\item В чем отличие k-Means от сети Кохонена (SOM)?
\begin{itemize}[label=$\square$, leftmargin=1.5cm, itemsep=0.2em, topsep=0.3em]
\item SOM учитывает топологию (соседство) кластеров, а k-Means ищет просто изолированные центры
\item k-Means обучает веса, а SOM нет
\item k-Means работает без учителя, а SOM - с учителем
\item Отличий нет
\end{itemize}
{\small\textit{Правильный ответ: SOM учитывает топологию (соседство) кластеров, а k-Means ищет просто изолированные центры.}}

\item Каким образом математически задается топологическое расстояние между двумя узлами на 2D-решетке сети Кохонена?
\begin{itemize}[label=$\square$, leftmargin=1.5cm, itemsep=0.2em, topsep=0.3em]
\item Расстоянием в исходном многомерном пространстве признаков
\item Расстоянием между их координатами в сетке (i, j) по Евклиду или Манхэттену
\item Квадратом ошибки
\item Скалярным произведением
\end{itemize}
{\small\textit{Правильный ответ: Расстоянием между их координатами в сетке (i, j) по Евклиду или Манхэттену.}}

\item Что означает экспоненциальное затухание радиуса окрестности в сети Кохонена?
\begin{itemize}[label=$\square$, leftmargin=1.5cm, itemsep=0.2em, topsep=0.3em]
\item Увеличение скорости обучения
\item Экспоненциальное затухание радиуса окрестности со временем
\item Вычисление Евклидова расстояния
\item Обратное распространение ошибки
\end{itemize}
{\small\textit{Правильный ответ: Экспоненциальное затухание радиуса окрестности со временем.}}

\item Почему нормализация входных данных критически важна для сети Кохонена?
\begin{itemize}[label=$\square$, leftmargin=1.5cm, itemsep=0.2em, topsep=0.3em]
\item Потому что алгоритм работает с Евклидовыми расстояниями, и признак с большим диапазоном значений будет несправедливо доминировать при поиске BMU
\item Чтобы уместить данные в оперативную память
\item Потому что алгоритм не поддерживает большие числа
\item Чтобы не было отрицательных значений
\end{itemize}
{\small\textit{Правильный ответ: Потому что алгоритм работает с Евклидовыми расстояниями, и признак с большим диапазоном значений будет несправедливо доминировать при поиске BMU.}}

\item Можно ли использовать обученную сеть Кохонена для классификации (когда метки классов становятся известны позже)?
\begin{itemize}[label=$\square$, leftmargin=1.5cm, itemsep=0.2em, topsep=0.3em]
\item Строго нет, алгоритм забудет данные
\item Да, если каждому узлу карты присвоить метку, определенную путем голосования по большинству попавших в него примеров из обучающей выборки
\item Да, но только если переучить сеть с Backpropagation
\item Да, сменив функцию активации
\end{itemize}
{\small\textit{Правильный ответ: Да, если каждому узлу карты присвоить метку, определенную путем голосования по большинству попавших в него примеров из обучающей выборки.}}

\end{enumerate}
